\chapter{\uppercase{Literature Review}}

\section{\uppercase{Traditional Static Analysis Approaches}}
\doublespacing{
Static Application Security Testing (SAST) has been the cornerstone of automated vulnerability detection for over two decades. Traditional SAST tools operate by analyzing source code or bytecode without executing the program, employing pattern matching, data flow analysis, and control flow analysis to identify potential vulnerabilities.
}

\subsection{Pattern-Based SAST Tools}
\doublespacing{
Early SAST tools such as Fortify\cite{fortify}, Checkmarx\cite{checkmarx}, and Coverity\cite{coverity} rely primarily on predefined rule sets and pattern matching. These tools scan code for known dangerous patterns such as unsanitized user input passed to SQL queries or system calls. While computationally efficient (typically $<$100ms per file), they suffer from high false positive rates (30-50\%) due to limited context understanding. Research by Johnson et al. (2013)\cite{johnson2013} demonstrated that developers ignore 40-60\% of SAST warnings due to alert fatigue from false positives.
}

\subsection{Taint Analysis}
\doublespacing{
Taint analysis, pioneered by Arzt et al. (2014)\cite{flowdroid} with FlowDroid, tracks the flow of untrusted data (taint sources) to security-sensitive operations (taint sinks). FlowDroid achieved significant advances in Android app security by implementing context-sensitive, flow-sensitive, object-sensitive, and lifecycle-aware taint tracking. However, FlowDroid's analysis is limited to intra-procedural flows and achieves only 29-40\% token-level coverage, struggling with implicit flows and field-sensitive tracking gaps that this project addresses through enhanced taint analysis with 22 custom propagation rules.
}

\subsection{CodeQL and Semantic Code Analysis}
\doublespacing{
GitHub's CodeQL\cite{codeql} represents a paradigm shift in static analysis by treating code as queryable data. CodeQL creates a relational database from source code, enabling analysts to write declarative queries in QL (a logic programming language) to identify complex vulnerability patterns. CodeQL's global taint tracking can follow data flow across entire codebases, making it significantly more powerful than traditional SAST. This project leverages CodeQL's infrastructure for CPG extraction, building upon its strengths while adding comprehensive taint analysis enhancements.
}

\section{\uppercase{Code Property Graphs}}

\subsection{CPG Foundation}
\doublespacing{
Yamaguchi et al. (2014)\cite{yamaguchi2014} introduced Code Property Graphs (CPGs) heterogeneous graphs combining Abstract Syntax Trees (AST), Control Flow Graphs (CFG), and Program Dependence Graphs (PDG). CPGs provide a unified representation capturing syntactic structure, control flow, and data dependencies. This unified view enables more sophisticated analysis than any individual representation alone.

The CPG framework models code as a graph $G = (V, E)$ where vertices $V$ represent program elements (statements, expressions, variables, functions) and edges $E$ represent relationships (control flow, data flow, syntactic containment). This graph-based representation enables graph traversal queries that can identify complex vulnerability patterns spanning multiple program constructs.
}

\subsection{Applications in Vulnerability Detection}
\doublespacing{
CPGs have been applied to various vulnerability detection tasks including:
\begin{itemize}
    \item \textbf{Missing security checks}: Identifying locations where security validation should occur but is absent
    \item \textbf{API misuse}: Detecting incorrect usage patterns of security-critical APIs
    \item \textbf{Information flow violations}: Tracking sensitive data propagation to unauthorized sinks
    \item \textbf{Anomaly detection}: Finding code patterns that deviate from established secure coding practices
\end{itemize}

This project extends the CPG framework with enhanced taint analysis capabilities, introducing four new edge types (TAINT\_FLOW, IMPLICIT\_FLOW, FIELD\_ACCESS, ALIAS\_OF) to capture comprehensive information flow relationships.
}

\section{\uppercase{Taint Analysis Techniques}}

\subsection{Classical Taint Tracking}
\doublespacing{
Classical taint tracking maintains a set of tainted variables and propagates taint through explicit data flow operations (assignments, function calls, return values). The analysis typically defines:
\begin{itemize}
    \item \textbf{Sources}: Points where untrusted data enters the program (user input, network, files)
    \item \textbf{Sinks}: Security-sensitive operations (SQL execution, command execution, file access)
    \item \textbf{Sanitizers}: Functions that remove taint (validation, encoding, escaping)
\end{itemize}

A vulnerability is reported when taint flows from a source to a sink without passing through an appropriate sanitizer. However, classical approaches miss many vulnerability classes due to limited propagation rules.
}

\subsection{Advanced Taint Propagation}
\doublespacing{
Recent work has extended taint analysis with:

\textbf{Implicit Flows}: Control-dependent information leaks where tainted data influences control flow, indirectly leaking information through branch decisions. For example, \texttt{if secret: x = 1} leaks information about \texttt{secret} through the value of \texttt{x}.

\textbf{Field-Sensitive Analysis}: Tracking taint through object fields and container elements independently, recognizing that \texttt{obj.field1} may be tainted while \texttt{obj.field2} is not.

\textbf{Alias Analysis}: Propagating taint through variable aliasing (\texttt{a = b; c = a}) and pointer relationships, essential for tracking taint in programs with complex data structures.

This project implements all three advanced techniques, achieving comprehensive coverage that surpasses existing tools.
}

\section{\uppercase{Research Gaps}}
\doublespacing{
Despite significant progress, several gaps remain in vulnerability detection research:

\begin{enumerate}
    \item \textbf{Limited Taint Coverage}: Existing taint analysis tools achieve only 29-40\% token-level coverage, missing implicit flows, field-sensitive propagation, and alias tracking.

    \item \textbf{High False Positive Rates}: Traditional SAST tools generate 30-50\% false positives, causing developer alert fatigue and tool abandonment.

    \item \textbf{Lack of Comprehensive CPG Infrastructure}: Most CPG research focuses on specific analysis tasks rather than building reusable, well-tested infrastructure for general-purpose program analysis.

    \item \textbf{Poor Explainability}: Existing tools provide limited explanation for detected vulnerabilities, reducing developer trust and hindering remediation.

    \item \textbf{Missing Extensibility}: Monolithic SAST architectures are difficult to extend with advanced techniques such as machine learning models or semantic analyzers.
\end{enumerate}

This project addresses gaps \#1, \#3, and \#5 by building comprehensive CPG infrastructure with enhanced taint analysis and modular architecture enabling future integration of advanced reasoning components (GNN, LLM, uncertainty quantification).
}

\section{\uppercase{Summary of Literature Review}}
\doublespacing{
Table \ref{tab:lit_review_summary} provides a comprehensive summary of the surveyed literature organized by research area.
}

\begin{table}[H]
\centering
\caption{Summary of Literature Review}
\label{tab:lit_review_summary}
\small
\renewcommand{\arraystretch}{1.2}
\begin{tabular}{|p{4cm}|p{3cm}|p{6cm}|}
\hline
\textbf{Research Area} & \textbf{Key Papers} & \textbf{Contributions} \\
\hline
\textbf{Traditional SAST} &
Johnson et al. (2013)\cite{johnson2013} &
Developer study showing 40--60\% of SAST warnings ignored due to false positives \\
\hline
\textbf{Commercial Tools} &
Fortify\cite{fortify}, Checkmarx\cite{checkmarx}, Coverity\cite{coverity} &
Enterprise SAST tools with pattern matching and data flow analysis; 30--50\% FP rates \\
\hline
\textbf{Taint Analysis} &
Arzt et al. (2014)\cite{flowdroid} &
FlowDroid: Context-, flow-, and object-sensitive taint tracking for Android; 29--40\% token coverage \\
\hline
\textbf{Query-Based Analysis} &
GitHub CodeQL\cite{codeql} &
Semantic code querying using QL language; global taint tracking; 29--35\% coverage \\
\hline
\textbf{Code Property Graphs} &
Yamaguchi et al. (2014)\cite{yamaguchi2014} &
Unified graph representation combining AST, CFG, and PDG for vulnerability detection \\
\hline
\textbf{Open-Source SAST} &
Semgrep\cite{semgrep} &
Lightweight pattern-based analysis; community rules; 20--25\% coverage, 40--50\% FP rate \\
\hline
\textbf{Benchmark Datasets} &
NIST Juliet\cite{juliet}, NVD\cite{nvd2023} &
Synthetic test cases (Juliet) and real-world CVEs (NVD) for evaluation \\
\hline
\end{tabular}
\end{table}
